%!TEX root = ../template.tex
%% ====================================================================
%% Capitulo 4 -- Materiais e metodos
%% Primeira versao completa 2026-09-08. Linguagem descritiva, sem rotulos
%% internos de organizacao do trabalho.
%% Fontes da vistoria (2026-09-08): ~/tese-extraccao (extraccao integral do
%% modelo Power BI v22 de 2026-07-31: modelo.bim, 159 medidas DAX, 21 queries M,
%% 32 particoes, metadados de reflexao, 25 TSV exportados, inventario e
%% manifesto MD5 das fontes, 33 prints), notas do Mind Palace (mapeamento e
%% validacao de jul-2026, memoria do projeto), PROTOCOL.md / CLAUDE.md e
%% docs/ do repositorio baselines-cc.
%% Marcas %% CONFIRMAR assinalam factos a confirmar pelo autor.
%% ====================================================================

\chapter{Materiais e métodos}
\label{cha:metodos}
% [REV-2026-09-13][C4-01][P2][ESTRUTURA]
\revisaotese{C4-01}{P2}{ESTRUTURA}
{Separar método reproduzível da história do desenvolvimento}
{Acrescentar uma secção inicial com desenho do estudo, unidade de análise,
janelas, perguntas e métricas primárias. Trazer para aqui as definições
dos candidatos, critérios e simulador atualmente introduzidos no Cap. 7.
Mover inventários extensos, erros de implementação, contagens de testes
e pormenores de publicação de corridas para apêndices. No corpo,
conservar decisões que alteram o estimando ou a validade dos resultados.}%
% [/REV-2026-09-13]

Este capítulo descreve o que foi construído e como foi avaliado. As três
primeiras secções tratam da plataforma: de onde vêm os dados, como foram
reconciliados num modelo dimensional único e que métricas se calculam sobre
ele. As duas últimas tratam do método de avaliação das \emph{baselines} que
os Capítulos~\ref{cha:diagnostico} e~\ref{cha:discussao} aplicam --- o
protocolo estatístico e a infraestrutura de reprodutibilidade que garante
que cada número citado nesta dissertação pode ser refeito. A matemática de
cada instrumento consta do Anexo~\ref{app:matematica}; aqui descreve-se o
desenho e as decisões que o justificam.

\section{Fontes de dados e ingestão}
\label{sec:fontes}

O consumo energético de uma unidade processual não existe numa fonte
única. Na refinaria de Sines, a informação necessária para o calcular está
repartida por três famílias de origens, com resoluções, unidades e
responsáveis diferentes, e a primeira decisão de desenho da plataforma foi
não escolher uma delas: lê-las todas e reconciliá-las por regras
declaradas.

% [REV-2026-09-13][C4-02][P1][DADOS]
\revisaotese{C4-02}{P1}{DADOS}
{Agregação diária e conversão energética precisam de equações}
{Distinguir caudal instantâneo/médio (t/h), totalizador (t) e energia.
Explicar fuso horário, dias incompletos, mínimo de horas válidas e regra
nos dias de 23/25 horas. Média x 24 só é justificável sob a convenção
temporal adotada. Para combustível com PCI variável, verificar se
energia = soma(massa\_h x PCI\_h) ou massa\_diária x PCI\_médio:
o segundo omite a covariância entre caudal e PCI. Quantificar o erro
ou documentar por que a aproximação é aceitável.}%
% [/REV-2026-09-13]
\paragraph{O historiador de processo.} O historiador IP21 regista as
medições de campo --- caudais de combustível e de vapor, carga das
unidades, densidades, temperaturas e pressões de \emph{flash}, poderes
caloríficos --- com resolução horária. O acesso faz-se através da
plataforma de dados da refinaria (Dremio), numa vista de negócio horária
por \emph{tag} (\texttt{ip21\_hourly}). A plataforma não consulta o
historiador \emph{tag} a \emph{tag}: gera, a partir da tabela de mapeamento
(Secção~\ref{sec:modelo-dimensional}), a instrução SQL que agrega cada
\emph{tag} ao dia com a regra que lhe cabe --- soma diária para caudais,
média diária para poderes caloríficos e propriedades, média multiplicada
por 24 para caudais horários, máximo para indicadores --- e devolve o
resultado em formato longo (\emph{tag}, ano, mês, dia, valor). A primeira
versão gerava uma coluna por \emph{tag} numa única consulta larga e
esgotava o tempo limite da plataforma; o formato longo, com uma consulta
por regra de agregação, resolveu-o sem perder nenhuma \emph{tag}.

\paragraph{Os balanços mensais.} O Processamento de Dados da refinaria
mantém dois livros de balanço mensal, com valores diários por referência
de medidor: o balanço de produção (BAL\_PROD; folhas de medidores e de
caudais, em toneladas por mês na coluna mensal e por dia nas colunas
diárias) e o balanço de utilidades (BAL\_UTIL; medidores em toneladas por
hora, vapor por nível de pressão). São a fonte oficial do fecho mensal e a
referência contra a qual os relatórios de consumo são emitidos. A leitura
é feita por funções de ingestão que percorrem as pastas de rede por ano e
mês e extraem, por referência, a série diária. Os dois livros não têm a
mesma estrutura: no de produção o primeiro dia do mês está seis colunas à
direita do identificador, no de utilidades quatro --- e a primeira versão
da ingestão, que assumia o desvio do primeiro em ambos, perdia dois dias
por mês no vapor e somava zeros nos dois últimos, sub-reportando os três
níveis de vapor em 5 a 7\%. O erro foi encontrado na validação da
Secção~\ref{sec:validacao-dados} e é o exemplo do que essa validação existe
para apanhar. O balanço de produção fornece ainda a folha de poderes
caloríficos diários dos combustíveis sem \emph{tag} própria (\emph{off-gas},
\emph{tail gas}, gás natural, coque), lida com uma guarda de gama física
(30 a 30\,000~kcal/kg, porque uma etiqueta repetida na folha produzia, sem
guarda, um poder calorífico médio de dois milhões) e com propagação do
último valor conhecido quando um mês não o traz --- registando a origem de
cada valor (\emph{bal}, \emph{locf}, \emph{backfill}) para que medido e
propagado nunca fiquem indistinguíveis.

\paragraph{A energia elétrica.} A eletricidade não tem série diária: existe
nos relatórios mensais de envio de cada fábrica, em MWh por unidade, com
blocos de consumo e de produção. A ingestão lê esses relatórios com um
filtro de nomes canónicos (padrão da fábrica seguido de espaço, sufixo de
envio, e uma guarda anti-duplicado por ano, mês e padrão que fica com o
ficheiro mais recente), porque um ficheiro de segurança deixado na pasta
de leitura --- uma versão anterior à revisão de janeiro de 2026 --- foi
somado ao definitivo durante várias atualizações, duplicando o consumo de
quatro unidades. Desde julho de 2026 a ingestão expõe o consumo bruto ao
lado do líquido (consumo menos produção): o líquido serve o reporte, o
bruto serve a \emph{baseline} e o índice de intensidade, porque uma unidade
que produz o que consome não deve aparecer com consumo nulo numa relação
com a carga.

\paragraph{Planos de fabrico, densidades e crudes.} Os planos de fabrico
semanais (folhas Excel no arquivo de planeamento, 2020 em diante)
fornecem o \emph{slate} de crudes previsto por período e por unidade; as
densidades de rede e as propriedades da carga (densidade, temperatura e
pressão de \emph{flash}, rendimento em resíduo) vêm do historiador e das
folhas de densidades, com um valor por omissão declarado no mapeamento
quando a \emph{tag} não existe --- e uma coluna de origem que diz, dia a
dia, se o valor foi medido ou assumido.

% [REV-2026-09-13][C4-03][P2][CONSISTENCIA]
\revisaotese{C4-03}{P2}{CONSISTENCIA}
{O inventário não fecha aritmeticamente}
{São anunciadas sete classes e 3 862 ficheiros, mas enumeram-se seis
classes que somam 3 809 (81+84+708+434+426+2 076): faltam identificar
53 ficheiros. Confirmar no manifesto sem corrigir por inferência.
Explicar se as seis divergências de hash afetam dados usados e qual
versão ficou congelada. Relacionar 32 tabelas do modelo, 25 exportadas
e 13 usadas na análise numa tabela compacta de âmbito/proveniência.}%
% [/REV-2026-09-13]
\paragraph{Inventário e extração.} Para congelar o estado das fontes e
tornar a plataforma reconstruível fora da rede da refinaria, fez-se em
julho de 2026 um inventário completo das origens (cerca de 67 mil ficheiros
percorridos, 14~GB) e uma extração das que o modelo consome: 3\,862
ficheiros e 988~MB em sete classes --- balanço de produção (81 ficheiros
mensais, cobertura completa de novembro de 2019 a julho de 2026), balanço
de utilidades (84), eletricidade (708), crudes (434), densidades (426),
planos de fabrico (2\,076) --- com o MD5 de cada ficheiro num manifesto e
uma verificação posterior (3\,856 iguais, 6 divergências, todas em
ficheiros de planos de fabrico de julho de 2024 alterados na origem entre a
cópia e a verificação).
O modelo do \emph{dashboard} foi serializado por inteiro na mesma data
(definição tabular, 159 medidas, 21 consultas de transformação, 32
partições, metadados de reflexão) e as 25 tabelas do modelo foram
exportadas em texto. É sobre essa exportação --- e não sobre uma ligação
viva --- que o Capítulo~\ref{cha:diagnostico} trabalha.

\section{Modelo dimensional e reconciliação}
\label{sec:modelo-dimensional}
% [REV-2026-09-13][F06][P2][FIGURA]
\revisaotese{F06}{P2}{FIGURA}
{Arquitetura de dados com granularidade e proveniência}
{Depois de apresentar as tabelas de factos, inserir um esquema com dois
painéis: (a) historiador/balanços/relatórios -\textgreater{} agregação e conversão -\textgreater{}
precedência por linha -\textgreater{} factos; (b) estrela simplificada energia/carga/
propriedades com calendário, unidade e vetor. Escrever a granularidade
de cada facto e marcar eletricidade mensal alocada. Não mostrar as 32
tabelas nem um screenshot ilegível do modelo inteiro. Referenciar o
dicionário e mostrar onde se retêm origem, flags e versões.}%
% [/REV-2026-09-13]

A plataforma é um modelo tabular em esquema em estrela~\cite{KimballRoss2013}, em modo de
importação, com 32 tabelas, 243 colunas, 13 relações e 159 medidas.
Três tabelas de factos concentram a informação: energia
(\texttt{FactEnergia\_Unified}: 200\,924 linhas dia$\times$unidade$\times$vetor,
de 2020-01-01 a 2026-07-31, 16 unidades, 7 vetores), carga
(\texttt{FactCarga\_Unified}: 45\,408 linhas desde novembro de 2019, com a
carga total e os seus componentes por nível de agregação) e propriedades
da carga (19\,105 linhas, 8 unidades), a que se juntam densidades, planos
de fabrico e crudes. As dimensões são o calendário (gerado, até ao dia
anterior à atualização), as unidades (16, em três fábricas mais os
extras, com o limiar de carga mínima de cada uma e a sua composição em
secções), os vetores (fuel gás, vapor a 24, 10 e 3~bar, energia
elétrica, coque e outros combustíveis, mais o agregado global), os crudes
(56) e duas dimensões de apresentação --- janela temporal e métrica ---
que alimentam os seletores do \emph{dashboard}. Os parâmetros físicos que
qualquer cálculo usa vivem em duas tabelas de parâmetros e nunca no
código: o poder calorífico de referência do gás natural
(11\,820~kcal/kg), as entalpias equivalentes dos três níveis de vapor
(698, 738 e 758~kcal/kg), a conversão da eletricidade (860,4~Mcal por MWh e um
factor de eficiência de conversão de 0,57), o número mínimo de
observações para estimar uma \emph{baseline} (30) e, por unidade, o limiar
de carga abaixo do qual um dia não é operacional --- 18\,600~t/d na
destilação atmosférica.

\paragraph{A tabela de mapeamento é o contrato.} O objeto central do
modelo não é um facto, é o mapeamento mestre: 286 linhas, cada uma a
ligar uma unidade física, um vetor e um papel (consumo ou produção) a
uma origem --- uma \emph{tag} do historiador, uma referência de balanço
ou as duas --- com o factor de unidade do canal de \emph{tags}, um factor
de inclusão do canal de balanço (1, 0, $-1$ ou uma fração de alocação),
a \emph{tag} de poder calorífico quando aplicável, datas de vigência e um
comentário de qualidade. Das 286 linhas, 235 apontam para o balanço de
utilidades e 51 para o de produção; 137 têm também \emph{tag} do
historiador; 213 são consumos e 73 produções. Cinco mapeamentos auxiliares
(poderes caloríficos, eletricidade, carga por \emph{tags} e por balanço,
densidades, propriedades) seguem o mesmo princípio: nenhuma constante de
conversão, nenhuma referência de medidor e nenhuma regra de alocação
existe fora de uma tabela versionada. Foi esta separação que permitiu
corrigir, em julho de 2026, a alocação da produção de vapor entre as duas
unidades de vácuo (24\% e 76\% a 24~bar; 69\% e 31\% a 3~bar, fórmulas
oficiais do controlo processual) alterando uma coluna e nenhuma linha de
código.

% [REV-2026-09-13][C4-04][P1][METROLOGIA]
\revisaotese{C4-04}{P1}{METROLOGIA}
{Consistência documental, reconciliação e exatidão física}
{Definir "reconciliação": aqui parece significar harmonização de fontes,
precedência e confronto, não reconciliação estatística sujeita a balanços
de massa/energia. Não inferir exatidão metrológica da concordância.
Documentar medidores a montante, calibração/qualidade disponíveis e
dependências entre canais. Acrescentar cobertura e discrepância por
ano/vetor para excluir mudanças de fonte/conversão como explicação da
deriva; cinco meses de validação não validam automaticamente seis anos.}%
% [/REV-2026-09-13]
\paragraph{Três canais, uma regra de precedência.} A energia de cada
dia, unidade e vetor chega por até três canais --- \emph{tags} do
historiador, balanço de produção ou utilidades, e relatórios de
eletricidade --- e a tabela unificada guarda os três lado a lado antes de
escolher. A regra é declarada e por linha: balanço, depois eletricidade,
depois \emph{tags}; nunca se mistura um canal com outro na mesma linha, e
uma coluna de origem regista qual venceu. Na carga a regra é a mesma
(balanço antes de historiador) e os componentes da carga seguem a fonte
vencedora. Nos 200\,924 registos de energia, 121\,828 têm balanço e
\emph{tags} em simultâneo --- é a redundância que a validação da
Secção~\ref{sec:validacao-dados} explora ---, 42\,132 só balanço e 36\,964
vêm dos relatórios de eletricidade. Uma quarta coluna, o consumo ``sem
produção'', é igual ao unificado em todos os vetores exceto na
eletricidade, onde usa o bruto: é a coluna que a \emph{baseline} e o índice
de intensidade consomem. Os relatórios operacionais e o reporte de
emissões continuam no líquido.

% [REV-2026-09-13][C4-05][P1][DADOS]
\revisaotese{C4-05}{P1}{DADOS}
{Imputação, zeros e seleção na variável resposta}
{"Não imputa" contradiz a propagação/backfill de PCI e os valores por
omissão de propriedades; restringir a afirmação ao consumo/carga sem
canal disponível. Quantificar imputações por variável/ano e a sua presença
no conjunto avaliado. Distinguir zero físico, dado ausente e zero fabricado,
sobretudo no vapor líquido. Excluir todos os zeros por resposta pode
alterar a população; mostrar sensibilidade e códigos de razão.
Backfill usa informação futura: declarar efeitos numa avaliação temporal.}%
% [/REV-2026-09-13]
\paragraph{O que o modelo não faz.} Não imputa: um dia sem valor em nenhum
canal fica sem valor, e um valor nulo ou zero do canal de balanço é
excluído da estimação de \emph{baselines} enquanto a origem dos zeros
fabricados não for resolvida na fonte. Não inclui o gás natural nem o
\emph{tail gas} no perímetro de nenhuma unidade, por decisão de âmbito.
E trata a eletricidade como o que ela é: uma alocação mensal distribuída
uniformemente pelos dias, que o Capítulo~\ref{cha:diagnostico} mostra não
ser uma medição diária.

\section{Cálculo de métricas energéticas}
\label{sec:metricas}

% [REV-2026-09-13][C4-06][P2][ENERGIA]
\revisaotese{C4-06}{P2}{ENERGIA}
{Escrever o balanço de conversão e justificar constantes}
{Incluir equações e tabela unidades de entrada -\textgreater{} fator -\textgreater{} t GNE/d.
Associar explicitamente 698/738/758 kcal/kg a cada pressão e indicar
origem, condições, referência entálpica, retorno de condensado e se são
equivalentes contabilísticos. Declarar PCI ou PCS de 11 820 kcal/kg.
Explicar eletricidade: E\_eq = MWh x 860,4/(11 820 x 0,57), com unidades
coerentes, e o significado do rendimento. Examinar sensibilidade dos
indicadores a constantes fixas. Evitar chamar eficiência termodinâmica
ao que depende de fronteiras e convenções de alocação.}%
% [/REV-2026-09-13]
\paragraph{Uma unidade comum.} Todos os vetores são convertidos a
toneladas de gás natural equivalente por dia (t~GNE/d), a unidade em que
a refinaria orçamenta e reporta. Os combustíveis (fuel gás, coque, outros)
convertem-se pelo caudal diário multiplicado pelo poder calorífico do dia
--- lido de uma \emph{tag} horária no fuel gás e no coque, da folha de
poderes caloríficos nos restantes --- dividido pelo poder calorífico de
referência do gás natural. O vapor converte-se pelo caudal multiplicado
pela entalpia equivalente do seu nível de pressão; o de 10~bar é um fluxo
líquido (consumo menos produção) e pode ser negativo. A eletricidade
converte-se de MWh para energia térmica e desta para combustível
equivalente pelo factor de eficiência de conversão. Cada regra está
declarada na tabela de parâmetros por vetor, e o cálculo é feito uma vez,
na ingestão, para que todas as medidas leiam a mesma coluna.

\paragraph{Dia operacional e consumo específico.} Um dia é operacional
para uma unidade quando a sua carga é igual ou superior ao limiar declarado
para essa unidade. O consumo específico é o total de GNE nos dias
operacionais dividido pela carga nesses mesmos dias, em t~GNE por
quilotonelada; as janelas móveis de 7 e 30 dias e o acumulado do período
são calculados sobre a mesma seleção. O limiar não é um parâmetro de
apresentação: é a fronteira do domínio em que a \emph{baseline} é estimada,
e a Secção~\ref{sec:perfil-dados} mostra o que ele exclui.

% [REV-2026-09-13][C4-07][P1][CRONOLOGIA]
\revisaotese{C4-07}{P1}{CRONOLOGIA}
{Precisar o que acontece efetivamente em janeiro}
{O método descrito estima cada ano com todos os dias disponíveis e
permite selecionar o ano de referência. Isso não demonstra que os alarmes
operacionais troquem automaticamente em janeiro para a reta do ano
anterior. Distinguir ajuste retrospetivo, ano corrente parcial e regra
real de aplicação. Afirmar mascaramento anual observado exige documentar
o histórico das referências usadas; sem ele, apresentar o mecanismo
como consequência possível do desenho/experiência de simulação.}%
% [/REV-2026-09-13]
\paragraph{A \emph{baseline} em produção.} Para cada unidade, vetor e ano
civil, o modelo estima por mínimos quadrados a relação
$E = a\,Q + b$ (o modelo estatístico de \emph{baseline} que a ISO~50006
tipifica~\cite{ISO50006_2014}) entre o consumo diário (em t~GNE/d, na coluna sem
produção) e a carga diária (em kton/d), sobre os dias operacionais do ano
com energia não nula e anteriores ao dia corrente. A estimação é feita
dentro do próprio modelo, numa tabela calculada que materializa, por
célula unidade$\times$vetor$\times$ano, o declive, a ordenada na origem,
o desvio-padrão residual, o coeficiente de determinação e o número de
observações --- 583 células no estado de julho de 2026, incluindo o ano
corrente marcado como parcial. O vetor global de cada unidade é a soma
dos vetores orçamentáveis (o coque, subproduto queimado no regenerador,
mantém reta própria mas fica fora da reta global). Foi esta tabela ---
exportada tal como o modelo a calcula --- que serviu de referência ao gate
de replicação da Secção~\ref{sec:replicacao}: o que o
Capítulo~\ref{cha:diagnostico} critica é, com igualdade numérica
verificada, o que este parágrafo descreve.

% [REV-2026-09-13][C4-08][P2][INDICADORES]
\revisaotese{C4-08}{P2}{INDICADORES}
{Excesso positivo não estima desperdício ou poupança}
{Definir e\_t = E\_t - previsão e distinguir soma(e\_t) de soma(max(e\_t,0)).
A segunda é positiva mesmo com erros centrados e cresce com ruído e
número de dias; não é desperdício, oportunidade de poupança nem critério
de eficiência entre vetores. Mostrar também desvio líquido, exposição
e escala. A soma simples de resíduos é um acumulado descritivo; reservar
CUSUM de deteção para a recorrência com referência k, limiar h e
calibração de sinais. Especificar lateralidade e contagem de episódios.}%
% [/REV-2026-09-13]
\paragraph{Resíduos, bandas e acumulados.} O resíduo diário é a diferença
entre o consumo observado e o previsto pela \emph{baseline} do ano de
referência selecionado, aplicada à carga do próprio dia; o resíduo
padronizado divide-o pelo desvio-padrão residual dessa \emph{baseline}. As
bandas a $\pm1$, $\pm2$ e $\pm3\sigma$ são as do desenho de produção, e a
regra de alarme é a saída da banda de $\pm2\sigma$. Sobre a série de
resíduos calculam-se o excesso diário (a parte positiva do resíduo), o
excesso acumulado por vetor, unidade e refinaria, a fração de dias acima
e abaixo da banda, o desvio médio fora da banda, a soma acumulada dos
resíduos desde o início do período selecionado (o gráfico de CUSUM~\cite{Page1954} do
\emph{dashboard}) e o ritmo de excesso nas janelas de 7 e 30 dias. Um
painel de regras de carta de controlo verifica oito regras clássicas de
causas especiais, na linha das de Nelson~\cite{Nelson1984}, sobre os resíduos padronizados --- um ponto fora de $\pm3\sigma$, dois em
três fora de $\pm2\sigma$, quatro em cinco fora de $\pm1\sigma$, oito
consecutivos do mesmo lado, seis monotónicos, catorze alternados, quinze
consecutivos dentro de $\pm1\sigma$, oito consecutivos fora de
$\pm1\sigma$ --- e conta as ocorrências de cada uma. Estas estatísticas
descrevem o método em produção; a sua validade estatística é a pergunta
do Capítulo~\ref{cha:diagnostico}.

% [REV-2026-09-13][C4-09][P2][INDICADORES]
\revisaotese{C4-09}{P2}{INDICADORES}
{Delimitar o EII e as emissões derivados da plataforma}
{Indicar versão/fonte autorizada do método Solomon, convenções de bruto/
líquido e se a aplicação diária é índice adaptado ou equivalente ao
benchmark oficial. Validar separadamente se alegar equivalência. Para
CO2, indicar fonte/unidade dos fatores e separar emissões diretas de
combustíveis e emissões atribuídas a vapor/eletricidade; um fator único
por GNE não demonstra emissões físicas por vetor. Se for proxy de
reporte interno, dizê-lo, sem ampliar o âmbito da tese.}%
% [/REV-2026-09-13]
\paragraph{O índice de intensidade energética.} Além da \emph{baseline}
interna, o modelo calcula por unidade o índice de intensidade energética
da metodologia de \emph{benchmarking} da indústria (Solomon): a razão,
em percentagem, entre o consumo real em kBTU e uma energia padrão que
depende da carga em barris e de um factor por unidade --- para a
destilação atmosférica, uma função da densidade da carga, do rendimento em
resíduo e das condições de \emph{flash}, com um teto para cargas leves ---
calculada só nos dias operacionais e com todas as propriedades disponíveis.
O índice é publicado no dia, em janelas móveis de 7 e 30 dias e no
período, contra as bandas de quartil da unidade, e decomposto por vetor
numa cascata. Serve de contexto externo; não é uma \emph{baseline} e não
entra na avaliação do Capítulo~\ref{cha:diagnostico}.

\paragraph{Emissões.} O equivalente em CO$_2$ é calculado do consumo
líquido por vetor com o factor de emissão declarado na dimensão dos
vetores, no dia e acumulado, e apresentado ao lado do GNE em todas as
vistas de síntese.
%% CONFIRMAR: unidade e origem do fator_co2 (DimVector, 2,5) — t CO2 por t GNE?

\section{Protocolo de avaliação estatística das \emph{baselines}}
\label{sec:protocolo-baselines}

\paragraph{A lição que definiu o protocolo.} Uma primeira análise das
\emph{baselines}, feita em abril de 2026 fora da plataforma, foi
invalidada em agosto: usava um limiar de carga escrito à mão (15~kton/d)
em vez do limiar oficial lido da configuração (18\,600~t/d), e por isso
analisava uma regressão que não era a que está em produção --- e concluía
coisas que não existem nos dados oficiais. Duas regras saíram daí e não
se negociaram depois: nenhum parâmetro de seleção é escrito à mão, todos
são lidos das tabelas de configuração do modelo de produção; e replica-se
antes de criticar --- se os coeficientes recalculados não reproduzem os de
produção em todas as células, o protocolo para. Uma terceira regra veio da
revisão externa do protocolo: não se criticam violações de independência
usando testes que assumem independência e apresentando os seus
$p$-values como prova.

\paragraph{Pré-registo e revisão adversarial.} O protocolo foi escrito
antes de qualquer execução, no sentido do pré-registo~\cite{Nosek2018} --- perguntas, estimadores, comparadores,
famílias de hipóteses, correção de multiplicidade, limiares e a grelha de
sensibilidade --- e submetido a quatro revisões adversariais independentes
(Anexo~\ref{app:adversarial}), que apontaram vinte defeitos na primeira
versão; a versão 1.2, que se descreve aqui, é a que resultou da
adjudicação desses defeitos. Todos os parâmetros vivem num ficheiro de
configuração versionado (\texttt{params.yaml}); nenhum é duplicado em
código. As poucas decisões tomadas depois de ver resultados --- sobre desenho,
nunca sobre critérios --- estão datadas no plano de execução e assinaladas
no texto onde entram.

\paragraph{Uma pergunta de dados antes de qualquer análise.} O primeiro
passo não é estatístico: é um perfil dos dados --- cobertura, mecanismo de
falta, distribuições brutas, efeito do filtro de seleção, séries temporais,
relação bivariada, consistência entre fontes --- que uma pessoa lê e
aceita, com declaração assinada e ligada por \emph{hash} à corrida que o
produziu, antes de qualquer instrumento arrancar. O erro de abril foi um
erro de seleção de dados que um perfil decente teria exposto; o gate
existe para que isso não volte a acontecer, e o código não tem caminho de
escrita para a declaração --- só uma pessoa a pode criar.

\paragraph{Replicação exata como gate bloqueante.} Sobre a exportação
congelada do modelo, o protocolo reestima a regressão anual de cada
unidade--vetor com o limiar lido da configuração e compara declive,
ordenada na origem, desvio-padrão residual e coeficiente de determinação
com os da tabela de \emph{baselines} exportada, célula a célula, com
tolerância pré-registada de $10^{-6}$; a identidade dos dias selecionados
é comparada por \emph{fingerprint} do conjunto. Trinta células têm de
reproduzir-se; se uma falha, nada a jusante corre. Cada instrumento
posterior repete uma guarda de replicação sobre as células que usa.

% [REV-2026-09-13][C4-10][P2][INFERENCIA]
\revisaotese{C4-10}{P2}{INFERENCIA}
{Definir população, alvo inferencial e limites de HAC/MBB}
{Dar tabela por teste com H0/estimando, unidade, pares/dias, comparador,
IC, família de multiplicidade e estatuto. HAC/MBB não tornam qualquer
série não-estacionária automaticamente tratável: discutir condições da
série de perdas/indicadores reamostrada, quebras e dependência persistente.
Precisar se os IC condicionam nos modelos já ajustados ou incorporam
reestimação do treino. Explicar blocos civis com lacunas (elegibilidade,
concatenação e reinício), não apenas citar o comprimento.}%
% [/REV-2026-09-13]
% [REV-2026-09-13][C4-11][P1][VALIDACAO]
\revisaotese{C4-11}{P1}{VALIDACAO}
{O alvo nominal de 95\% não foi validado pelo teste a 90\%}
{357/400 = 89,25\% de cobertura simulada. Não rejeitar cobertura inferior
a 90\% não valida um IC nominal de 95\%; o alvo foi ainda escolhido após
conhecer o resultado. Manter este ensaio como diagnóstico e declarar
prévia ou posteriormente os critérios com precisão. Mostrar cobertura
por rho, n e bloco com IC Monte Carlo; não transferir validação do bloco
de 60 dias para o bloco base de sete. Resultados fronteiriços precisam
da qualificação no local, além da secção final de limitações.}%
% [/REV-2026-09-13]
\paragraph{As perguntas e a inferência.} As conclusões do diagnóstico
assentam em cinco perguntas, cada uma com estimador, comparador e
critério declarados: se o modelo de um ano prevê o seguinte melhor do que
referências triviais (média do ano de treino como nulo primário; mesmo dia
do ano anterior e média móvel de trinta dias como secundários~\cite{HyndmanKoehler2006}), medido
pela diferença emparelhada de perda absoluta diária~\cite{DieboldMariano1995,HarveyLeybourneNewbold1997}; se os declives anuais
são iguais, por um teste de Wald sobre uma regressão empilhada com
interações ano$\times$carga, na linha do teste de Chow~\cite{Chow1960}; se a banda de $\pm2\sigma$ do ano de treino
cobre a fração nominal dos dias do ano seguinte; que desvio mínimo o
sistema de alarmes deteta e em quanto tempo; e se covariáveis
pré-especificadas (densidade da carga, temperatura e pressão de
\emph{flash}) acrescentam valor preditivo fora da amostra. A inferência
respeita a dependência serial que os próprios resíduos exibem: erros-padrão
robustos à heterocedasticidade e à autocorrelação~\cite{NeweyWest1987}
(\emph{kernel} quadrático-espectral com largura de banda de
Andrews~\cite{Andrews1991}) para os testes de igualdade, e intervalos percentil~\cite{Efron1979,EfronTibshirani1993} por
\emph{bootstrap} de blocos móveis~\cite{Kunsch1989,Lahiri2003} --- bloco
base de $\lceil n^{1/3}\rceil$ dias~\cite{HallHorowitzJing1995} com
desfasamentos civis exatos, 2\,000 réplicas, semente fixa, sensibilidade a blocos de 7, 14,
30 e 60 dias --- para todos os contrastes de perdas e coberturas. Uma
conclusão existe apenas quando o intervalo exclui o valor de referência;
nunca por comparação de pontos. A multiplicidade é corrigida dentro da
família dos testes de igualdade de declives (Benjamini--Hochberg~\cite{BenjaminiHochberg1995}, com
Benjamini--Yekutieli~\cite{BenjaminiYekutieli2001} como verificação); os restantes resultados são por
contraste e são apresentados como tal. A validação do \emph{bootstrap} por
simulação, a $\rho = 0{,}90$ e bloco de 60 dias, deu uma cobertura de
357/400; o resultado foi adjudicado por teste binomial exato~\cite{ClopperPearson1934}
depois de conhecido e fica declarado como o que é --- ausência de evidência de
cobertura inferior a 0,90, não prova de cobertura nominal.

\paragraph{Diagnósticos descritivos.} Um segundo conjunto de instrumentos
--- dependência e variância dos resíduos, distribuição e caudas,
segmentação do espaço operacional, sobreposição de suporte entre anos,
influência de semanas, plausibilidade física dos coeficientes, correlação
entre vetores --- documenta mecanismos. O seu estatuto é declarado
inferior: os $p$-values são índices descritivos calculados sob
independência, nenhum rótulo é conclusão, e nenhuma afirmação da
dissertação assenta neles.

\paragraph{Sensibilidade às convenções.} Todas as conclusões são
reavaliadas numa grelha pré-registada que cruza o limiar de carga a 0,8,
1,0 e 1,2 vezes o oficial com duas partições temporais --- anos civis e
janelas de doze meses ancoradas a 1 de julho --- num total de 624
reavaliações, com a célula de referência a reproduzir exatamente os testes
deterministas. A grelha qualifica; nunca cria conclusões.

% [REV-2026-09-13][F07][P2][FIGURA]
\revisaotese{F07}{P2}{FIGURA}
{Linha temporal dos três desenhos de avaliação}
{Inserir aqui, antes da demonstração: três faixas alinhadas em 2020--2026.
(a) treino y -\textgreater{} teste y+1 nos cinco pares; (b) mesmos pares para comparar
candidatos, com afinação interna; (c) referência 2020 -\textgreater{} aplicação
2021--2025. Marcar 2026 parcial e indicar que os dados já foram explorados.
Uma faixa auxiliar distingue observação real de simulação. Legenda
declara disponibilidade de Q\_t/covariáveis e dos resíduos passados,
para separar normalização ex post de previsão genuína a um passo.}%
% [/REV-2026-09-13]
\paragraph{Demonstração das linhas orientadoras.} A avaliação de
candidatos alternativos (estimadores robustos, bandas com dependência
modelada, critério de aceitação) usa os mesmos cinco pares anuais, com
constantes fixadas a priori ou afinadas apenas no ano de treino, e
conclusões só por intervalo. A formulação original previa uma partição
selada com 2025 intocado; foi substituída, por decisão datada, por
validação temporal encadeada~\cite{Tashman2000,BergmeirHyndmanKoo2018},
porque 2025 já tinha informado o diagnóstico --- e nenhum resultado dessa parte se chama \emph{holdout}. O
critério de aceitação (conteúdo preditivo e banda suportada, cada um em
pelo menos quatro dos cinco pares) foi fixado antes de qualquer resultado.

\paragraph{Referência congelada e simulação.} A exploração da
Secção~\ref{sec:referencia-congelada} congela a relação de 2020, a sua
escala (por validação cruzada em seis blocos contíguos~\cite{BergmeirHyndmanKoo2018})
e o seu estado
inicial, aplica-os a 2021--2025 numa população primária declarada (dias
retidos, dentro do suporte de carga de 2020, com observação e predição
finitas), compara modelos apenas na interseção dessas populações, e mede o
comportamento de \emph{baselines} e monitores com verdade conhecida num
gerador declarado (relação de 2020 com as covariáveis de referência, ruído
autorregressivo, máscara real de dias admissíveis, limiares calibrados sob
a hipótese nula a 5\%, quatrocentas réplicas por célula, oito contrastes
co-primários declarados antes da execução). Todos os artefactos
distinguem, linha a linha, o que foi observado no \emph{snapshot} do que
foi imposto em simulação.

\section{Reprodutibilidade e proveniência}
\label{sec:reprodutibilidade}
% [REV-2026-09-13][C4-12][P2][REPRODUTIBILIDADE]
\revisaotese{C4-12}{P2}{REPRODUTIBILIDADE}
{Integridade não certifica correção científica}
{Hashes e corridas imutáveis são pontos fortes, mas não provam que
filtros, estimandos ou interpretação estejam certos. Substituir "provar
que não foi editado" por "detetar divergências face a um manifesto
identificado", com a âncora de confiança. No corpo, apresentar um exemplo
número -\textgreater{} tabela -\textgreater{} corrida -\textgreater{} dados/versão; transferir detalhes de
infraestrutura para apêndice. Prever pacote público sanitizado com código,
ambiente e dados sintéticos, se os dados industriais não forem divulgáveis.}%
% [/REV-2026-09-13]

A exigência que estrutura toda a infraestrutura é simples de enunciar:
qualquer número citado nesta dissertação tem de poder ser refeito, e tem
de ser possível provar que não foi editado depois de
produzido~\cite{Peng2011,Sandve2013}.

\paragraph{Dados congelados e identificados.} As 13 tabelas do modelo
dimensional que a avaliação consome foram exportadas em 2026-07-31 e
copiadas para o repositório de análise como ficheiros de texto só de
leitura; o \emph{hash} SHA-256 de cada um está registado num ficheiro de
fontes, e uma corrida cujo \emph{hash} de entrada não coincida falha em vez
de produzir números diferentes em silêncio. A identidade das observações
selecionadas pela \emph{baseline} de produção é guardada como
\emph{fingerprints} de referência.

\paragraph{Corridas write-once com manifesto.} Cada execução de um
instrumento é uma corrida: prepara-se numa área de \emph{staging},
verifica-se antes de começar (dados, ambiente, árvore de código limpa) e
publica-se por renomeação atómica para uma pasta identificada por data e
\emph{hash}, ou para uma pasta de falhas se a verificação final não
passar. Nunca se edita uma corrida publicada; uma corrida substituída
fica publicada e é apenas deixada de citar. O manifesto de cada corrida
regista os \emph{hashes} dos dados de entrada, o \emph{commit} do código, a
exigência de que a árvore esteja limpa, as versões exatas do ambiente
(dois ambientes fixados por ficheiro de bloqueio, um para a avaliação do
método de produção e outro para a exploração com referência congelada, que
não podem ser trocados), as sementes, e o inventário de todos os ficheiros
produzidos com o seu SHA-256 e um \emph{hash} do próprio inventário. As
interfaces de execução não aceitam parâmetros científicos: só a raiz do
repositório e a pasta de saída; tudo o resto é lido da configuração.

\paragraph{Pins e verificação independente.} Os resultados que a
dissertação cita são corridas canónicas identificadas (``pinadas'') na
configuração; qualquer instrumento que consuma resultados de outro lê-os
da corrida pinada, verifica o \emph{hash} do manifesto de origem antes de
ler e recusa o lote errado. Um verificador que não importa nada do
repositório --- só as bibliotecas de \emph{hash} e de leitura de texto ---
recomputa, para cada corrida citada, os \emph{hashes} e tamanhos de todos
os ficheiros e o \emph{hash} do inventário. À data desta versão, 110
manifestos e 1\,481 ficheiros foram verificados sem divergência por uma
auditoria independente.

\paragraph{Números nos documentos.} Os relatórios de resultados não
contêm números escritos à mão: cada valor da exploração com referência
congelada é derivado das tabelas pinadas por um gerador de documentos e
registado num bloco de proveniência no fim de cada documento, que um teste
automático recalcula e confronta com a prosa. As tabelas dos capítulos de
resultados desta dissertação seguem a mesma origem; a Secção~\ref{sec:limitacoes}
diz onde a garantia é mais fraca do que o mecanismo sugere.

\paragraph{Testes e ambiente.} O código de avaliação tem 1\,061 testes na
suite do método de produção e 683 na da exploração com referência
congelada, cada um no seu ambiente; sem teste, um instrumento não está
feito, e um teste que verifica a presença de um campo não verifica o
significado do filtro que o produziu --- lição registada depois de uma
revisão ter apanhado cinco afirmações que a tabela por trás delas refutava.
O determinismo é declarado com precisão: as corridas são reprodutíveis no
ambiente fixado com igualdade numérica até cerca de $10^{-15}$; a igualdade
byte a byte de ficheiros com valores em vírgula flutuante depende da
biblioteca de álgebra linear da máquina, e não é prometida.

\paragraph{Trilha adversarial.} As revisões externas do protocolo, as
auditorias científicas de cada parte (a integridade da cadeia de
resultados foi aprovada e a conformidade científica reprovada em setembro
de 2026, com correções feitas no código e testes que reproduzem cada
achado, e uma segunda reinspeção ainda em curso à data desta versão) e as
decisões do autor que delas resultaram estão versionadas no repositório e
resumidas no Anexo~\ref{app:adversarial}. É a parte da metodologia que
mais custou e a que mais vale: cinco afirmações da primeira versão dos
resultados eram falsas, todas do mesmo tipo --- narrativa à frente da
verificação ---, e nenhuma sobreviveu a abrir a linha da tabela que a
devia sustentar.
